% Limitations

\section{Limitations}
\label{sec:limitations}

We state the limitations of the present work explicitly, since several of
them are deliberate scope choices rather than oversights.

\begin{itemize}
  \item \textbf{Two dimensions, periodic box, single forcing mode.}
  All validation is for the 2D incompressible Navier--Stokes equations on
  the periodic box $\T^2$ with the single-mode forcing of
  eq.~\eqref{eq:forcing} (Section~\ref{sec:problem}). Free-slip walls,
  boundary layers, and multi-mode or spatially heterogeneous forcing are
  outside the scope of this work; the stream-function representation and the
  exact viscous step both rely on the periodic setting and do not carry over
  to bounded domains without additional treatment.

  \item \textbf{Reynolds numbers up to 5000.} The turbulent-validation
  ladder of Section~\ref{sec:setup} stops at $\mathrm{Re} = 5000$. Whether
  the adaptive rank and the accuracy of Section~\ref{sec:res-error} hold
  above that Re is untested; the cost model of Section~\ref{sec:cost}
  suggests the per-step picture should be stable (the nonlinear residual
  remains $O(n \log n)$ and rank-independent), but this is an expectation,
  not a result.
  % [RESOLVED: the Re = 5000 runs show no degradation -- the retained rank is
  % the same at all three Reynolds numbers and the horizon varies by under 8%
  % across them (Section~\ref{sec:res-error}) -- so this stays a limitation
  % about Re > 5000 rather than a statement about the runs we have.]

  \item \textbf{Rank decay is only fully exercised in the laminar run.}
  The tolerance-based decay rule is exercised only where a state genuinely loses
  modes. In the Taylor--Green decay the state is a single Fourier mode from the
  outset, so its numerical rank is one and the criterion must hold it there
  rather than spend rank on a spectrum that is not present; the run does. Under
  sustained forcing the rank instead rises, and it does so to the budget, which bounds what the decay rule can be tested against
  here.

  \item \textbf{No per-step speedup claim.} Because the nonlinear residual
  is evaluated on the full grid at rank-independent $O(n \log n)$ cost
  (Section~\ref{sec:cost}), SP-DLRA is not expected to beat the full-grid
  spectral reference per step at small rank, and the cost section reports
  the regimes where it is \emph{slower} (Figure~\ref{fig:cost}). The
  contribution is the reduced-model capability with exact
  divergence-freeness and a rank criterion that is measured, and whose
  saturation at the grid's alias-free ceiling --- a rank budget we set from the
  grid, $2\lfloor N/3\rfloor+1$, not a rank the grid imposes or the dynamics chooses
  (Section~\ref{sec:res-error}) --- we report as a finding,
  validated in a high-Reynolds-number turbulent regime; a wall-clock
  advantage is
  neither claimed nor assumed.

  \item \textbf{No stability analysis of the full coupled scheme.} The
  viscous step is exact within the ansatz and the nonlinear step is a
  standard projected midpoint update, but the scheme with online rank
  changes has no stability or error analysis in this work. The residual
  indicator \eqref{eq:indicator} is a heuristic error control, not a proven
  a~priori bound; long-time behavior over the retained windows is supported
  by the runs themselves (Sections~\ref{sec:res-error}
  and~\ref{sec:res-fidelity}), not by a theorem.

  \item \textbf{Forcing-aware invariant not yet formalized for the reduced
  model.} The unforced energy monotonicity (I2) is replaced under forcing by
  the forcing-aware balance of eq.~\eqref{eq:energy}; the precise
  reduced-model statement that the runs are checked against is being
  defined by the theoretical-research line of work (decision D3) and is not
  yet finalized. We therefore report the energy balance of
  Section~\ref{sec:res-fidelity} as an empirical check pending that
  statement.
  % [PENDING-THEORETICAL-RESEARCH: replace this paragraph with the final
  % forcing-aware invariant statement once it lands.]

  \item \textbf{No statistical-convergence study.} The retained statistical
  windows of Section~\ref{sec:setup} are fixed in advance; we do not study
  how the statistics (kinetic-energy spectrum, dissipation rate) converge as
  the window length grows, nor do we run longer than the retained windows.
  The rank is already at its ceiling for $99\%$ of the longer run, so what
  longer runs leave open is the statistical quantities, not the rank.
\end{itemize}
